\section{结论与局限性}
\label{sec:conclusion}

我们发布 TwinRouterBench，一个建立在路由最重要场景之上的步级路由基准：包括 SWE-bench、PinchBench、BFCL、mtRAG 和 QMSum 在内的长上下文、多轮智能体工作负载。静态赛道通过贪心的顺序锁定降级搜索，近似每一步的最低成本充分层级；针对多轮轨迹的人工审计确认了所得标签的可靠性。在动态赛道上，100 样例的留出 SWE-bench 评测证明在线执行框架能产出有意义的端到端结果。一个在静态标签上训练的逻辑回归路由器，在动态赛道上以相同解决率实现 53\% 的成本降低，表明静态数据不仅是快速离线评测工具，也是面向实用路由器的有效训练监督。

\paragraph{局限性。}\label{sec:limitations}
来自 520 个实例的 970 个路由步骤构成的静态语料库，足以进行初步诊断和训练，但并未穷尽全部智能体工作负载、领域或路由模式。标签是在 \S\ref{sec:static-scoring}的固定模型池、级联协议和价格前沿下估计的；若未来模型同时变得低价且高能力，就需要更新标签。将动态覆盖扩展到 SWE-bench Verified 之外，并随模型池和定价演进保持数据最新，是未来工作。

\paragraph{更广泛影响。}
通过提供标准化的步级评测，TwinRouterBench 可能加速高成本效益 LLM 路由器的开发，间接降低智能体部署的 API 成本和能源消耗。双赛道设计也将确定性离线评分与在线验证分开，从而改善可复现性。一个潜在负面影响是基准过拟合：社区工作可能集中于当前的 970 条语料和五类工作负载，而牺牲未覆盖的领域、语言或智能体架构。此外，层级标签绑定于固定模型池和价格快照；随着模型与价格演进，陈旧标签可能误导路由器训练或评测。我们通过版本化模型池与定价、明确的范围限制，以及支持更新模型池和重新标注的发布设计来缓解这些风险。
